% Examen ALEA du 9 janvier 2026, exercice 4.
\uuid{RM9l}
\titre{Assurance vie, loi du nombre de décès et bénéfice annuel}
\niveau{L2}
\module{Probabilité}
\chapitre{Probabilité discrète}
\sousChapitre{Approximation de la loi binomiale}
\theme{Loi binomiale, approximation, assurance}
\auteur{}
\datecreate{2026-10-01}
\organisation{AMSCC}
\difficulte{}
\contenu{
\texte{Une compagnie assure $N=10\,000$ clients sur la vie, qui paient chacun une prime annuelle de $A$ euros. On estime que chaque client a une probabilité de décès au cours d'une année égale à $p=0{,}006$, indépendamment les uns des autres. La prime de décès est de $B$ euros.}
\begin{enumerate}
  \item \question{Quelle est la loi du nombre annuel de décès ? Comment peut-on l'approcher ?}
  \item \question{Si $B=1\,000$, pour quelles valeurs de $A$ la compagnie a-t-elle une probabilité strictement inférieure à $0{,}01$ d'être en déficit ?}
  \item \question{Si $A>15$, pour quelles valeurs de $B$ la compagnie a-t-elle une probabilité strictement supérieure à $0{,}7$ de faire un bénéfice annuel supérieur à $50\,000$ euros ?}
\end{enumerate}
}
